\chapter{Introduction}

The ability to operate a computer that is physically somewhere else has moved
from a specialist convenience to an assumed capability. A support engineer
expects to reach a customer's laptop without travelling to it; an administrator
expects to reach a rack-mounted server that has no monitor attached; a student
expects to reach the laboratory machine that holds a licensed toolchain. Each of
these situations is served today, but rarely by the same tool and rarely without
a compromise the user is asked to accept silently.

This chapter states the motivation for building a new remote desktop system,
defines the problem precisely, lists the objectives that guided the design, and
outlines how the remainder of the report is organised.

\section{Overview}

A remote desktop system performs three tasks continuously and must perform all
three well. It \emph{captures} the visual state of the controlled machine, it
\emph{transports} that state to the controlling machine, and it \emph{returns}
keyboard and pointer events in the opposite direction. The difficulty is that
these tasks pull against one another. Capturing at high fidelity produces more
data than a residential uplink can carry; compressing aggressively introduces
delay that makes the pointer feel detached from the hand; and transporting
anything at all requires that two machines which may both be hidden behind
network address translation somehow find a path to each other.

Existing systems resolve this tension in different places. Microsoft's Remote
Desktop Protocol achieves excellent efficiency by transmitting drawing
primitives rather than pixels, but it is tied to the Windows display stack and
assumes a routable address or a corporate gateway. Virtual Network Computing
transmits framebuffer rectangles and is admirably portable, but it was designed
for local networks and its bandwidth behaviour on a wide-area link is poor.
Commercial products such as TeamViewer and AnyDesk solve reachability elegantly
by routing every session through vendor-operated relays, and in doing so place
the vendor in a position of complete visibility over traffic the user may regard
as sensitive.

\updesk{} takes the position that the reachability problem and the trust problem
are separable. Reachability is a well-studied networking question with a
standard answer --- the ICE framework, with STUN for address discovery and TURN
as a fallback relay. Trust is a cryptographic question, and the answer is to
ensure that whatever infrastructure assists with reachability is structurally
incapable of reading the session. WebRTC provides both: ICE for connectivity and
mandatory DTLS-SRTP for media encryption, with keys negotiated directly between
the endpoints. A signaling server is still required to introduce the peers, but
it handles only metadata.

\section{Motivation}

Three observations motivated this work.

\textbf{Relay-mediated architectures concentrate risk.} When every frame of a
support session traverses a third party's servers, the confidentiality of that
session rests entirely on the operator's policy and competence rather than on
mathematics. A design in which the coordinating server never holds a media key
removes that dependence. This is not a hypothetical concern for an
administrator entering credentials on a remote machine during a support call.

\textbf{Unattended access is where most remote desktop tools quietly fail.} A
system that requires a human at the far end to click ``accept'' is adequate for
support but useless for the server in the corner that needs attention at two in
the morning. Genuine unattended operation requires the host to run before any
user logs in, to survive reboots, and to remain usable while Windows displays
its login screen or a User Account Control prompt. Those screens are drawn on a
separate, isolated desktop that an ordinary user-session application cannot
capture. Meeting this requirement is a systems problem, not a user-interface
problem, and it drove much of the architecture described in Chapter~4.

\textbf{Systems languages have become practical for this class of software.}
Remote desktop hosts handle untrusted network input while manipulating raw
framebuffers --- historically a rich source of memory-safety vulnerabilities.
Rust offers memory safety without a garbage collector, and its asynchronous
ecosystem is now mature enough to implement WebRTC, TLS termination and video
encoding in one coherent codebase. Building \updesk{} predominantly in Rust was
therefore both an engineering and a security decision.

\section{Problem Statement}

\emph{To design and implement a cross-platform remote desktop system in which
screen and control data travel over a directly negotiated, end-to-end encrypted
peer-to-peer channel; in which the coordinating server authenticates endpoints
cryptographically but cannot decrypt session content; and which additionally
provides a fully unattended host capable of operating from system boot,
including on privileged desktops that conventional user-session applications
cannot capture.}

\section{Objectives}

The work was directed at the following objectives.

\begin{enumerate}
  \item To design a signaling protocol that authenticates hosts and controllers
        using Ed25519 public-key cryptography, resists replay through a
        challenge--response exchange, and exposes no session content to the
        server.
  \item To implement that signaling server in Rust with optional TLS termination
        and optional persistence of enrolled identities and session audit
        records.
  \item To implement a graphical host application and a graphical controller
        application that establish a WebRTC peer connection, stream video in one
        direction and control events in the other.
  \item To implement a headless native host that captures the framebuffer
        without any user gesture, encodes it to H.264, and serves it over a
        native WebRTC peer connection.
  \item To implement native keyboard and pointer injection so that the native
        host supports control and not merely observation.
  \item To package the native host as a Windows service so that it is present
        from boot, survives reboot and session switching, and can reach the
        active desktop.
  \item To evaluate the resulting system for connection-establishment success
        across NAT configurations, interactive latency, frame rate and bandwidth
        consumption.
\end{enumerate}

\section{Scope and Limitations}

The implementation targets Windows as the host platform, because the unattended
and secure-desktop requirements are inherently platform-specific and Windows
presents the hardest version of that problem. The controller application is
portable, and exploratory Android clients exist in the repository but are not
claimed as completed contributions. Audio forwarding and file transfer were
deliberately excluded in order to keep the video and control paths correct
before adding surface area. Multi-monitor capture is limited to the primary
display. These boundaries are revisited in Chapter~7 as directions for further
work.

\section{Organisation of the Report}

Chapter~2 surveys existing remote desktop protocols and products and the
underlying real-time communication technologies, and identifies the gap this
work addresses. Chapter~3 records the functional, non-functional and
environmental requirements. Chapter~4 presents the system architecture: the
component decomposition, the signaling protocol, the connection-establishment
sequence, and the security model. Chapter~5 describes the implementation of each
subsystem with reference to the actual codebase. Chapter~6 reports the
experimental setup and the measurements obtained. Chapter~7 concludes and
proposes future work.
